\section{Conclusion}

The research question of this chapter asked how gamification can be effectively applied to a gardening system to keep users more engaged in a normal activity. Based on the theory and the Plant Up! implementation, the answer is that gamification works best when it does not replace the activity, but gives the activity clearer feedback, visible progress, and a stronger emotional structure.

For Plant Up!, this means that the app should not reward users simply for opening screens or repeating actions. The important point is that game elements have to stay connected to real plant care. Quests, experience points, levels, streaks, rewards, and social features are useful when they guide users toward checking their plant, understanding its condition, and reacting only when care is actually needed. In this way, gamification turns plant care from a passive routine into an active process without making the plant itself less important.

The chapter also showed that effective gamification needs balance. Rewards can motivate users, but they should not become the only reason to use the app. Streaks can build habits, but they must not encourage harmful behavior such as watering every day. Community features can increase motivation, but they should support learning and shared progress instead of creating unfair pressure. This is especially important in Plant Up!, because every plant has different needs and every user starts from a different situation.

In conclusion, gamification can keep users more engaged in a gardening system when it is designed as a support layer for meaningful care behavior. Plant Up! applies this by connecting digital progress to sensor-based feedback and plant-specific recommendations. The result is an app concept where motivation, feedback, and responsibility work together: the user stays engaged not because the app distracts from plant care, but because it makes plant care easier to understand, easier to follow, and more rewarding over time.
